\documentclass[../../../include/open-logic-section]{subfiles}

\begin{document}

\olfileid{sth}{choice}{justifications}
\olsection{Panimbangan Intrinsik Babagan Pilihan}

Pitakonan kang luwih amba yaiku apa Pengurutan Becik, Pilihan, utawa
kabandingan ukurané kabèh himpunan---endi waé, amarga padha
ekuivalen---bisa diwènèhi pambenaran.

Ing ngisor iki ana upaya pambenaran kang \emph{intrinsik}. Ing
\olref[z][story]{sec}, kita ngenalaké sawatara prinsip babagan
hirarki. Salah sijiné prayoga dipratelakaké maneh:
\begin{enumerate}
	\item[] \stagesacc. Kanggo sembarang tahap $S$, lan sembarang
	himpunan kang wis kawangun \emph{sadurungé} tahap $S$,
	ing tahap $S$ kawangun himpunan kang unsuré pas himpunan-himpunan
	kasebut. Ora ana himpunan liya kang kawangun ing tahap~$S$.
\end{enumerate}
Pancen, akèh panulis kang nyaranaké yèn Aksioma Pilihan bisa
dibeneraké liwat prinsip iki utawa prinsip kang sarupa. Kita bakal
mènèhi gambaran cekak bab pendekatan kasebut.

Kita miwiti saka asil prasaja kang mènèhi \emph{ekuivalensi liya}
tumrap Pilihan:

\begin{thm}[ing $\ZF$]\ollabel{choiceset}
Pilihan ekuivalen karo prinsip iki. Yèn !!{element}s saka $A$
pisah-pisah sepasang-sepasang lan ora kosong, mula ana $C$ kang
ndadèkaké $C \cap x$ singleton kanggo saben $x \in A$. ($C$ kaya
mangkono diarani {himpunan pilihan} kanggo $A$.)
\end{thm}

Bukti asil iki prasaja, lan kita tinggalaké minangka gladhen kanggo
pamaos.

\begin{prob}
Buktèkna \olref[sth][choice][justifications]{choiceset}. Yèn
panjenengan kangèlan, buktiné bisa ditemokaké ing
\cite[kaca~242--3]{Potter2004}.
\end{prob}

Pokoké, himpunan pilihan kanggo $A$ iku mung jangkauan sawijining
fungsi pilihan kanggo $A$. Mula kanggo mbeneraké Pilihan, kita bisa
nyoba mbeneraké pratelan kang ekuivalen iki liwat anané himpunan
pilihan. Saiki kita bakal nyoba nindakaké iku.

Anggepen !!{element}s saka $A$ pisah-pisah lan ora kosong.
Miturut \stageshier{} (delengen \olref[z][story]{sec}), $A$
kawangun ing sawatara tahap~$S$. Gatekna yèn kabèh !!{element}s
saka $\bigcup A$ wis ana sadurungé tahap $S$. Saiki, miturut
\stagesacc{}, kanggo \emph{sembarang} himpunan kang kawangun
sadurungé~$S$, kawangun himpunan kang unsuré pas himpunan-himpunan mau.
Kanthi tembung liya, saben kumpulan kang \emph{bisa} kawangun saka
himpunan kang wis ana sadurungé bakal ana ing tahap~$S$. Nanging
mesthi \emph{bisa} milih obyèk kang mbentuk sawijining himpunan
pilihan kanggo~$A$; iki mung sawijiné himpunan pérangan tartamtu
saka $\bigcup A$. Dadi ana himpunan pilihan kang dikarepaké.

Iki upaya kang \emph{cekak banget} kanggo mbeneraké Pilihan kanthi
alasan intrinsik. Nanging yèn panjenengan kepengin nyinaoni luwih
adoh, wacanen pangrembakan Potter kang becik
(\citeyear[\S14.8]{Potter2004}).

\end{document}
